% cap_resultados.tex
\chapter{Resultados e Discussão}
\label{chap:resultados}

Este capítulo apresenta os resultados do protocolo Sapo descrito no Capítulo \ref{chap:materiais_metodos}. A discussão separa quatro níveis de evidência: a comparação principal com hiperparâmetros selecionados no bloco de desenvolvimento, a estabilidade com múltiplas sementes dos modelos selecionados, as linhas de base ingênuas e as ablações usadas para interpretar decisões de modelagem.

As métricas usadas no capítulo seguem as definições da Seção \ref{sec:metricas_avaliacao}: erro médio absoluto (MAE), raiz do erro quadrático médio (RMSE), erro percentual médio absoluto (MAPE), coeficiente de determinação (\(R^2\)), erro médio de viés (MBE) e MAE em eventos com PM2,5 \(\geq 35\) (Peak35 MAE). Em todos os casos, a referência de avaliação é a medição real de PM2,5 disponível no teste.

\section{Caracterização da Comparação Principal}

A comparação principal usa a estação Sapo, tarefa 48 para 24 horas e divisão cronológica 70/15/15. A seleção de hiperparâmetros e ablações ocorre no bloco de desenvolvimento, e a avaliação é feita no bloco final de teste. O conjunto de dados final contém 34.991 registros horários. No conjunto consolidado, a proporção de valores originalmente ausentes é de aproximadamente 19,26\% para material particulado fino (PM2,5) e 19,58\% para material particulado com diâmetro de até 10 micrômetros (PM10). Esses percentuais reforçam a importância de manter indicadores de ausência e de avaliar os modelos contra medições reais disponíveis.

Os eventos altos são raros no teste. Na caracterização exploratória do bloco final, pontos observados com PM2,5 \(\geq 35\) representam cerca de 0,96\% do teste. Portanto, o MAE global é dominado por regimes baixos e moderados. A análise de cauda é apresentada separadamente para evitar que o melhor erro médio seja confundido com melhor resposta a picos.

\section{Resultado Principal no Protocolo Final}

A Tabela \ref{tab:resultado_principal_sapo_hpo} apresenta os resultados principais no bloco final de teste. Todos os modelos foram reexecutados no mesmo contrato final, com 48 horas de entrada, 24 horas de horizonte, conjunto \texttt{clean\_pm10\_decoder\_proxy}, mesma divisão cronológica e as mesmas medições de PM2,5 como referência. As LSTMs e o XGBoost usam hiperparâmetros selecionados por validação temporal \textit{walk-forward}; o novo Seq2Seq atenção usa a variante selecionada pelas ablações de arquitetura, perda e janela de entrada, reexecutada aqui no mesmo protocolo comparável. O melhor desempenho global na semente canônica foi obtido pela LSTM direta, enquanto o novo Seq2Seq atenção ficou em segundo por MAE e apresentou melhor comportamento em picos.

\begin{table}[htb]
    \centering
    \caption{Resultado principal no protocolo final de Sapo.}
    \label{tab:resultado_principal_sapo_hpo}
    \resizebox{\textwidth}{!}{%
    \begin{tabular}{lrrrrrr}
        \toprule
        \textbf{Modelo} & \textbf{MAE} & \textbf{RMSE} & \textbf{MAPE} & \textbf{\(R^2\)} & \textbf{H1 MAE} & \textbf{H24 MAE} \\
        \midrule
        LSTM direta & \textbf{2,7713} & \textbf{3,8346} & \textbf{23,01} & \textbf{0,5154} & 2,6725 & \textbf{2,8670} \\
        Seq2Seq atenção novo & 2,8132 & 3,8799 & 23,23 & 0,5039 & 2,6222 & 2,9138 \\
        XGBoost com saída múltipla & 2,8786 & 3,9793 & 23,80 & 0,4781 & \textbf{2,5618} & 3,0273 \\
        Seq2Seq atenção canônica & 2,8811 & 3,9859 & 23,79 & 0,4764 & 2,7155 & 2,9743 \\
        LSTM recursiva & 2,8876 & 3,8380 & 24,39 & 0,5146 & 2,7380 & 3,1492 \\
        Seq2Seq básico & 2,9553 & 4,0967 & 23,99 & 0,4469 & 2,6519 & 3,1514 \\
        \bottomrule
    \end{tabular}%
    }
\end{table}

O resultado reforça a conclusão central: a LSTM direta permanece vencedora por erro médio. O novo Seq2Seq atenção melhora substancialmente a posição da família Seq2Seq em relação à variante ponderada anterior, reduzindo o MAE da semente canônica de 2,8658 para 2,8132 e aproximando RMSE e \(R^2\) dos valores da LSTM direta. O XGBoost continua sendo uma referência forte baseada em árvores sobre atributos derivados das janelas temporais e manteve o menor erro no primeiro horizonte, mas perdeu para LSTM direta e novo Seq2Seq no MAE global.

Entre os modelos neurais, a diferença entre LSTM direta e novo Seq2Seq é pequena no erro médio, mas não desaparece: 2,7713 contra 2,8132 na semente canônica. A LSTM recursiva apresentou RMSE e \(R^2\) próximos aos da LSTM direta, mas teve pior MAE e maior erro no horizonte final. O Seq2Seq básico e o Seq2Seq atenção canônica ficaram abaixo da nova variante, mostrando que o ganho veio da combinação específica de desenho do decodificador, perda ponderada, agrupamento de atenção por defasagens e seleção da janela de 48 horas, não apenas do uso genérico de uma arquitetura codificador-decodificador.

\section{Estabilidade com Múltiplas Sementes}

A Tabela \ref{tab:multi_seed_sapo} resume a sensibilidade à semente aleatória para os modelos selecionados. A média inclui a semente canônica 42 e três sementes adicionais: 7, 21 e 123.

\begin{table}[htb]
    \centering
    \caption{Estabilidade com múltiplas sementes dos modelos selecionados.}
    \label{tab:multi_seed_sapo}
    \resizebox{\textwidth}{!}{%
    \begin{tabular}{lrrrrr}
        \toprule
        \textbf{Modelo} & \textbf{MAE médio} & \textbf{Desvio MAE} & \textbf{RMSE médio} & \textbf{\(R^2\) médio} & \textbf{MBE médio} \\
        \midrule
        LSTM direta & \textbf{2,8055} & 0,0540 & 3,8611 & 0,5086 & 0,0435 \\
        Seq2Seq atenção novo & 2,8218 & 0,0220 & 3,8989 & 0,4989 & 0,0651 \\
        LSTM recursiva & 2,8425 & 0,0426 & \textbf{3,8006} & \textbf{0,5239} & 0,4505 \\
        Seq2Seq atenção canônica & 2,8727 & 0,0175 & 3,9592 & 0,4834 & 0,1641 \\
        XGBoost com saída múltipla & 2,8852 & \textbf{0,0077} & 3,9895 & 0,4755 & -0,0344 \\
        Seq2Seq básico & 2,9444 & 0,0113 & 4,0778 & 0,4520 & -0,1777 \\
        \bottomrule
    \end{tabular}%
    }
\end{table}

A média com múltiplas sementes mantém a LSTM direta como melhor modelo por MAE, mas o novo Seq2Seq atenção passa a ser o segundo colocado com margem pequena: MAE médio 2,8218 contra 2,8055. A LSTM recursiva passa a ter melhor RMSE e \(R^2\) médios, ao custo de viés positivo consistente. Essa diferença confirma que a escolha do vencedor depende da métrica priorizada. Para o objetivo principal do TCC, que usa MAE como métrica central de erro médio no teste, a LSTM direta continua sendo a escolha principal; para eventos altos, a análise de cauda favorece o novo Seq2Seq.

\section{Linhas de Base Ingênuas}

A Tabela \ref{tab:baselines_ingenuos_sapo} apresenta linhas de base simples avaliadas no mesmo bloco final de teste. A pontuação de habilidade (\textit{skill score}) usa a persistência \(y(t)\) como referência, de forma que valores positivos indicam redução de MAE em relação a essa linha de base.

\begin{table}[htb]
    \centering
    \caption{Linhas de base ingênuas no bloco final de teste.}
    \label{tab:baselines_ingenuos_sapo}
    \resizebox{\textwidth}{!}{%
    \begin{tabular}{lrrrr}
        \toprule
        \textbf{Linha de base} & \textbf{MAE} & \textbf{RMSE} & \textbf{\(R^2\)} & \textbf{Habilidade MAE} \\
        \midrule
        Persistência \(y(t)\) & 3,8971 & 5,2834 & 0,0801 & 0,0000 \\
        Sazonal \(y(t-24)\) & 3,3739 & 4,7174 & 0,2666 & 0,1342 \\
        Sazonal \(y(t-48)\) & 3,7599 & 5,2441 & 0,0937 & 0,0352 \\
        Média móvel 6h & 3,3481 & 4,5533 & 0,3167 & 0,1409 \\
        Média móvel 24h & \textbf{3,0377} & \textbf{4,2234} & \textbf{0,4122} & \textbf{0,2205} \\
        \bottomrule
    \end{tabular}%
    }
\end{table}

A melhor linha de base ingênua foi a média móvel de 24 horas, com MAE 3,0377. A LSTM direta selecionada por \textit{walk-forward} reduz esse erro para 2,7713, uma melhora de aproximadamente 8,8\% sobre a melhor linha de base simples e de 28,9\% sobre a persistência.

\section{Erro por Horizonte}

A Tabela \ref{tab:erro_por_horizonte_sapo} detalha o MAE nos horizontes 1, 6, 12 e 24. O XGBoost teve o melhor primeiro horizonte, mas sua degradação até o horizonte 24 foi maior que a da LSTM direta. A LSTM direta apresentou evolução mais estável ao longo do horizonte, o que ajuda a explicar sua vantagem no MAE agregado.

\begin{table}[htb]
    \centering
    \caption{MAE por horizonte no teste.}
    \label{tab:erro_por_horizonte_sapo}
    \resizebox{\textwidth}{!}{%
    \begin{tabular}{lrrrrr}
        \toprule
        \textbf{Modelo} & \textbf{H1} & \textbf{H6} & \textbf{H12} & \textbf{H24} & \textbf{H24 - H1} \\
        \midrule
        LSTM direta & 2,6725 & \textbf{2,7218} & \textbf{2,7773} & \textbf{2,8670} & \textbf{0,1945} \\
        Seq2Seq atenção novo & 2,6222 & 2,7788 & 2,8268 & 2,9138 & 0,2916 \\
        XGBoost com saída múltipla & \textbf{2,5618} & 2,8004 & 2,9092 & 3,0273 & 0,4655 \\
        Seq2Seq atenção canônica & 2,7155 & 2,8572 & 2,8923 & 2,9743 & 0,2588 \\
        LSTM recursiva & 2,7380 & 2,7338 & 2,9061 & 3,1492 & 0,4112 \\
        Seq2Seq básico & 2,6519 & 2,9035 & 2,9850 & 3,1514 & 0,4995 \\
        \bottomrule
    \end{tabular}%
    }
\end{table}

A estabilidade por horizonte é relevante porque a tarefa é multi-horizonte. Um modelo pode ser atraente no primeiro horizonte e ainda perder desempenho agregado quando a janela avança. Neste experimento, o XGBoost manteve o menor erro em H1, mas teve degradação maior até H24. A LSTM direta combinou bom erro inicial com menor degradação e melhor H24. O novo Seq2Seq ficou entre esses extremos: não venceu H1 nem H24, mas manteve erro por horizonte competitivo e melhor que as variantes Seq2Seq anteriores.

\begin{figure}[htb]
    \centering
    \includegraphics[width=\linewidth]{erro_por_horizonte_sapo.png}
    \caption{Evolução do MAE por horizonte no bloco final de teste para os modelos com métricas consolidadas.}
    \label{fig:erro_por_horizonte_sapo}
\end{figure}

\section{Picos, Cauda e Suavização}

A Tabela \ref{tab:cauda_sapo} mostra métricas associadas à amplitude e aos regimes altos. O desvio-padrão observado no teste é aproximadamente 5,51, e o percentil 99 (p99) resume o valor abaixo do qual se concentram 99\% das previsões ou observações. Valores da razão entre desvios-padrão abaixo de 1 indicam que a série prevista é mais suave que a série observada.

\begin{table}[htb]
    \centering
    \caption{Métricas de amplitude e cauda no teste.}
    \label{tab:cauda_sapo}
    \resizebox{\textwidth}{!}{%
    \begin{tabular}{lrrrrr}
        \toprule
        \textbf{Modelo} & \textbf{Peak35 MAE} & \textbf{p99 previsto} & \textbf{p99 observado} & \textbf{Razão desvio-padrão} & \textbf{MBE} \\
        \midrule
        LSTM direta & 13,6112 & 28,57 & 34,00 & 0,7391 & -0,1366 \\
        Seq2Seq atenção novo & \textbf{11,0815} & \textbf{32,50} & 34,00 & \textbf{0,9002} & 0,0433 \\
        XGBoost com saída múltipla & 12,2583 & 30,73 & 34,00 & 0,8506 & 0,0066 \\
        Seq2Seq atenção canônica & 14,7027 & 29,13 & 34,00 & 0,8113 & 0,0709 \\
        LSTM recursiva & 12,0392 & 30,49 & 34,00 & 0,8751 & 0,4450 \\
        Seq2Seq básico & 12,9806 & 30,65 & 34,00 & 0,8342 & -0,3855 \\
        \bottomrule
    \end{tabular}%
    }
\end{table}

Essas métricas mostram um compromisso importante. A LSTM direta vence em erro médio, mas apresenta razão entre desvios-padrão de 0,7391 e p99 previsto de 28,57, abaixo do p99 observado de 34,00. Isso sugere uma previsão mais suave e com subestimação de parte dos eventos altos. O novo Seq2Seq atenção preserva mais amplitude, com p99 previsto de 32,50, razão de desvio-padrão 0,9002 e menor Peak35 MAE. Portanto, a escolha do modelo depende do objetivo: minimizar erro médio favorece a LSTM direta; priorizar resposta a eventos altos favorece o novo Seq2Seq. A Figura \ref{fig:predito_observado_pico_sapo} reapresenta uma janela do teste em painéis separados por modelo, com escala vertical comum e interrupções nos trechos sem horários observados. Essa organização evita que as curvas disputem o mesmo eixo visual e deixa explícito que as quebras na linha observada correspondem a lacunas da própria série de teste.

\begin{figure}[htb]
    \centering
    \includegraphics[width=\linewidth]{predito_observado_pico_sapo.png}
    \caption{Predições agregadas no teste em uma janela com evento de maior concentração, separadas por modelo e com interrupções nas lacunas temporais.}
    \label{fig:predito_observado_pico_sapo}
\end{figure}
\clearpage

\subsection{Análise por Evento de Pico}

Além das métricas ponto a ponto, os picos foram agregados por horário único e agrupados em eventos contínuos com PM2,5 observado \(\geq 35\). Foram encontrados 21 eventos no teste. A Tabela \ref{tab:eventos_pico_sapo} resume o erro médio dentro dos eventos, o viés médio nos eventos, a taxa de subestimação e o erro no horário do pico observado.

\begin{table}[htb]
    \centering
    \caption{Métricas por evento contínuo de pico no teste.}
    \label{tab:eventos_pico_sapo}
    \resizebox{\textwidth}{!}{%
    \begin{tabular}{lrrrr}
        \toprule
        \textbf{Modelo} & \textbf{MAE evento} & \textbf{MBE evento} & \textbf{Subestimação} & \textbf{Erro no pico observado} \\
        \midrule
        Seq2Seq atenção novo & \textbf{12,4138} & \textbf{-11,6337} & \textbf{0,8651} & \textbf{-12,7452} \\
        LSTM recursiva & 12,8186 & -12,8186 & 1,0000 & -13,6797 \\
        XGBoost com saída múltipla & 13,7253 & -13,7253 & 1,0000 & -14,7604 \\
        Persistência \(y(t)\) & 13,8803 & -13,8689 & 0,9894 & -14,9206 \\
        Sazonal \(y(t-24)\) & 14,3677 & -14,2487 & 0,9405 & -15,3809 \\
        Seq2Seq básico & 14,3725 & -14,3725 & 1,0000 & -15,5662 \\
        LSTM direta & 15,4683 & -15,4683 & 1,0000 & -16,3894 \\
        Seq2Seq atenção canônica & 15,8503 & -15,8503 & 1,0000 & -17,0143 \\
        \bottomrule
    \end{tabular}%
    }
\end{table}

A leitura por evento é mais severa que o MAE global: todos os modelos subestimam os picos, mas em graus diferentes. O novo Seq2Seq atenção tem o menor erro médio por evento, o menor viés negativo e a menor taxa de subestimação, enquanto a LSTM direta, apesar de vencer no erro global, é uma das mais conservadoras nos eventos altos. Essa evidência explicita o compromisso entre erro médio e resposta a episódios de maior concentração.

\section{Novo Seq2Seq com Atenção}

A atenção canônica com Huber não foi o melhor caminho durante o desenvolvimento, mas as ablações do Seq2Seq produziram uma variante mais forte para o contrato final. A nova configuração mantém o codificador-decodificador com atenção, usa perda ponderada por regime, remove o \textit{input feeding}, reinicia o estado do decodificador por horizonte, agrupa a atenção em \textit{buckets} de defasagem e adiciona uma perda auxiliar leve para eventos. A análise de janela indicou que 48 horas continuam sendo o melhor compromisso para essa variante: janelas de 72, 96 e 168 horas não melhoraram o MAE nem os picos.

\begin{table}[htb]
    \centering
    \caption{Novo Seq2Seq atenção no protocolo final.}
    \label{tab:novo_seq2seq}
    \resizebox{\textwidth}{!}{%
    \begin{tabular}{lrrrrr}
        \toprule
        \textbf{Variante} & \textbf{MAE} & \textbf{RMSE} & \textbf{\(R^2\)} & \textbf{Peak35 MAE} & \textbf{p99 previsto} \\
        \midrule
        Seed 42 & 2,8132 & 3,8799 & 0,5039 & 11,0815 & 32,50 \\
        Média com múltiplas sementes & 2,8218 & 3,8989 & 0,4989 & 11,5855 & 31,90 \\
        Desvio com múltiplas sementes & 0,0220 & 0,0574 & 0,0147 & 0,3771 & 0,56 \\
        \bottomrule
    \end{tabular}%
    }
\end{table}

O novo Seq2Seq substitui a variante antiga com \texttt{weighted\_l1} como representante principal da família Seq2Seq. Na semente canônica, o MAE caiu de 2,8658 para 2,8132 e o Peak35 MAE caiu de 12,3697 para 11,0815 em relação ao Seq2Seq ponderado anterior. Ainda assim, a LSTM direta permanece melhor em MAE global. A interpretação mais defensável é que o novo Seq2Seq não vence a comparação por erro médio, mas oferece o melhor compromisso entre erro agregado e preservação de picos.

\section{Ablação PM10 Causal}

A Tabela \ref{tab:ablation_pm10_causal} compara o contrato final contra duas remoções de informação: manter PM10 apenas no histórico, sem substitutos causais no decodificador, e remover PM10 completamente, deixando PM2,5 e atributos temporais. Os hiperparâmetros foram mantidos fixos a partir da seleção \textit{walk-forward}; essas ablações não fizeram nova otimização de hiperparâmetros (HPO) e foram avaliadas no mesmo bloco final de teste.

\begin{table}[htb]
    \centering
    \caption{Ablação de PM10 causal no contrato final.}
    \label{tab:ablation_pm10_causal}
    \resizebox{\textwidth}{!}{%
    \begin{tabular}{llrrrr}
        \toprule
        \textbf{Modelo} & \textbf{Contrato} & \textbf{MAE} & \textbf{\(\Delta\) MAE} & \textbf{RMSE} & \textbf{\(R^2\)} \\
        \midrule
        LSTM direta & PM10 causal completo & \textbf{2,7713} & 0,0000 & \textbf{3,8346} & \textbf{0,5154} \\
        LSTM direta & PM10 histórico sem substituto no decodificador & 2,8575 & +0,0862 & 3,9678 & 0,4812 \\
        LSTM direta & PM2,5 + tempo apenas & 2,8456 & +0,0743 & 3,9569 & 0,4840 \\
        XGBoost & PM10 causal completo & \textbf{2,8786} & 0,0000 & \textbf{3,9793} & \textbf{0,4781} \\
        XGBoost & PM10 histórico sem substituto no decodificador & 2,8873 & +0,0088 & 4,0063 & 0,4710 \\
        XGBoost & PM2,5 + tempo apenas & 2,8956 & +0,0170 & 4,0194 & 0,4676 \\
        \bottomrule
    \end{tabular}%
    }
\end{table}

A ablação confirma que os substitutos causais de PM10 foram uma decisão central, sobretudo para a LSTM direta: remover esses atributos aumentou o MAE em 0,0862, e remover PM10 completamente aumentou o MAE em 0,0743. No XGBoost, o ganho global foi menor, mas ainda consistente em MAE, RMSE e \(R^2\). Assim, o papel de PM10 causal deve ser apresentado como contribuição metodológica do contrato de dados, não apenas como detalhe de engenharia de atributos.

\section{Limite Informacional com Variáveis Futuras}

Após a ablação causal de PM10, foi executado um diagnóstico adicional com Seq2Seq atenção para responder a uma pergunta diferente: quanto do erro restante parece vir de informação futura ausente, e não apenas de arquitetura ou hiperparâmetros? Para isso, foi usado um cenário \textit{oracle future auxiliary}, no qual o decodificador recebe PM10 e partículas totais em suspensão (PTS) observados no próprio horizonte futuro. Esse cenário não representa um uso causal de previsão e não deve ser tratado como candidato a modelo final, mas funciona como limite superior informacional sob a mesma estação e tarefa.

Na estação Sapo, a tentativa de usar um conjunto mais amplo de variáveis mostrou uma limitação prática: entre as variáveis auxiliares avaliadas, PM10 e PTS foram as únicas com observações úteis de forma consistente; variáveis meteorológicas e gases estavam ausentes ou esparsos demais no recorte analisado. A Tabela \ref{tab:oracle_future_aux_sapo} resume os resultados diagnósticos.

\begin{table}[htb]
    \centering
    \caption{Diagnóstico \textit{oracle future auxiliary} no Seq2Seq atenção em Sapo.}
    \label{tab:oracle_future_aux_sapo}
    \resizebox{\textwidth}{!}{%
    \begin{tabular}{lrrrr}
        \toprule
        \textbf{Contrato diagnóstico} & \textbf{MAE} & \textbf{RMSE} & \textbf{\(R^2\)} & \textbf{Peak35 MAE} \\
        \midrule
        PM10 causal no decodificador & 2,7777 & 3,9131 & 0,4954 & 10,2406 \\
        PM10 + PTS causal no decodificador & 2,8758 & 3,9949 & 0,4741 & 10,5980 \\
        PM10 futuro real no decodificador (\textit{oracle}) & \textbf{2,4201} & 3,2659 & 0,6485 & 7,9835 \\
        PM10 + PTS futuro real no decodificador (\textit{oracle}) & 2,4223 & \textbf{3,2654} & \textbf{0,6486} & \textbf{7,1718} \\
        \bottomrule
    \end{tabular}%
    }
\end{table}

O ganho do \textit{oracle} é grande e ajuda a interpretar o problema. Em relação ao controle causal com PM10 e PTS, o uso de variáveis futuras reais reduziu o MAE de 2,8758 para 2,4223 e elevou o \(R^2\) de 0,4741 para 0,6486. A melhoria em picos também foi clara: o Peak35 MAE caiu de 10,5980 para 7,1718. Esse resultado mostra que existe informação relevante no futuro de variáveis auxiliares, especialmente em PM10, que o contrato causal não disponibiliza no instante da previsão.

Ao mesmo tempo, o próprio \textit{oracle} impõe uma leitura conservadora: mesmo com PM10 e PTS futuros reais, o \(R^2\) ficou em torno de 0,65, ainda distante de 1. Portanto, a distância até um \(R^2\) quase perfeito não parece ser apenas consequência de HPO insuficiente ou de uma arquitetura Seq2Seq subótima. Parte relevante da variância horária de PM2,5 em Sapo provavelmente depende de fatores não observados no contrato atual, como meteorologia local, estabilidade e mistura atmosférica, transporte espacial, emissões, ressuspensão e episódios súbitos.

\section{Seq2Seq Multi-Alvo com Variáveis Auxiliares}

Também foi testada uma adaptação do Seq2Seq atenção para prever PM2,5 como alvo principal e, durante o treinamento, prever PM10 e PTS futuros como tarefas auxiliares. A hipótese era que a aprendizagem conjunta poderia induzir uma representação interna mais sensível à dinâmica dos poluentes correlacionados, sem usar variáveis futuras reais no decodificador.

\begin{table}[htb]
    \centering
    \caption{Experimentos Seq2Seq multi-alvo com PM10 e PTS como tarefas auxiliares.}
    \label{tab:seq2seq_multialvo_sapo}
    \resizebox{\textwidth}{!}{%
    \begin{tabular}{lrrrr}
        \toprule
        \textbf{Variante diagnóstica} & \textbf{MAE} & \textbf{RMSE} & \textbf{\(R^2\)} & \textbf{Peak35 MAE} \\
        \midrule
        Seq2Seq causal PM10 + PTS, sem multi-alvo & 2,8758 & 3,9949 & 0,4741 & \textbf{10,5980} \\
        Multi-alvo, peso auxiliar 0,03 & 2,7997 & 4,0045 & 0,4715 & 14,5285 \\
        Multi-alvo, peso auxiliar 0,08 & 2,8876 & 4,1483 & 0,4329 & 10,5677 \\
        Multi-alvo, dimensão 128 e peso auxiliar 0,12 & \textbf{2,7689} & \textbf{3,9373} & \textbf{0,4891} & 11,3975 \\
        \bottomrule
    \end{tabular}%
    }
\end{table}

O melhor multi-alvo reduziu o MAE em relação ao controle causal com PM10 e PTS, de 2,8758 para 2,7689, e elevou o \(R^2\) de 0,4741 para 0,4891. Esse é um indício positivo de que a tarefa auxiliar pode regularizar a representação do Seq2Seq. Porém, o ganho não reproduziu o salto observado no \textit{oracle} e ainda piorou a métrica de pico em relação ao controle causal. Em termos práticos, aprender PM10 e PTS futuros não equivale a conhecer PM10 e PTS futuros.

Esses dois diagnósticos convergem para a mesma interpretação. A arquitetura pode extrair algum ganho adicional, mas o principal gargalo para \(R^2\) muito alto está na informação disponível. Sem variáveis meteorológicas e espaciais observadas de forma confiável, e sem acesso causal ao futuro dos poluentes auxiliares, o modelo tende a suavizar a variância horária e subestimar parte dos picos.

\section{Interpretação da Atenção}

Os diagnósticos anteriores dos pesos de atenção foram mantidos como evidência auxiliar, mas não são usados como argumento principal para a nova comparação final. Nas variantes avaliadas, os pesos médios de atenção não formaram um seletor temporal concentrado nem uma explicação direta da previsão. Essa cautela é coerente com a crítica de \citeonline{jain2019attention}: pesos de atenção podem ajudar a inspecionar o modelo, mas não devem ser tratados automaticamente como explicação completa da decisão.

Na rodada final, o melhor comportamento do novo Seq2Seq aparece de forma mais clara nas métricas de erro, amplitude e eventos do que na interpretação isolada da atenção. Portanto, a leitura adotada no TCC é pragmática: a atenção faz parte da arquitetura vencedora dentro da família Seq2Seq, mas o ganho observado parece depender da combinação entre contrato causal de variáveis, janela de 48 horas, desenho do decodificador, perda ponderada e seleção por ablações. Não há evidência suficiente para afirmar que o modelo aprendeu uma regra temporal simples baseada apenas em uma defasagem específica.

\section{Ablações Auxiliares e Resultados Negativos}

As ablações internas do Seq2Seq tiveram papel importante no desenvolvimento do protocolo, mas nem todas pertencem ao mesmo escopo comparável da tabela principal. Estratégias como Luong \textit{dot/general/local}, atenção temporal de múltiplas cabeças (\textit{multi-head}), portões (\textit{gates}), atenção por grupos, atenção guiada (\textit{guided attention}), DA-RNN, viés de recência (\textit{recency bias}) e variantes de \textit{scheduled sampling} foram úteis para testar hipóteses, mas não produziram ganho determinístico suficiente para alterar o vencedor global atual.

O principal achado dessas ablações não é que atenção seja irrelevante em qualquer cenário, mas que, neste conjunto de dados e sob este contrato, mudanças no contrato de dados, no uso causal de PM10 e na função de perda tiveram efeito mais claro que modificações internas sofisticadas no mecanismo de atenção. O resultado negativo também é informativo: arquiteturas mais complexas não compensam automaticamente problemas de dados, suavização ou objetivo de treinamento desalinhado com picos.

A comparação fixa inicial Sapo 70/15/15, executada sem otimização de hiperparâmetros (HPO) e com outro conjunto de variáveis, deve ser lida como etapa histórica. Ela ajudou a organizar a comparação inicial entre quatro modelos de aprendizado profundo e XGBoost, mas não é a fonte da conclusão final. A conclusão atual se baseia no protocolo final com PM10 causal, HPO \textit{walk-forward}, avaliação com múltiplas sementes dos modelos selecionados e teste no mesmo conjunto final.

\section{Ameaças à Validade}

Há limitações importantes. Primeiro, o eixo principal usa uma estação, Sapo, o que limita a generalização para outras estações e regiões. Segundo, o \textit{walk-forward} foi usado para HPO dentro dos 85\% anteriores ao teste, mas o teste final continua sendo um único bloco cronológico reservado; isso preserva o cenário temporal, mas ainda não equivale a repetir todo o estudo em múltiplas estações ou períodos independentes. Terceiro, os modelos neurais apresentam sensibilidade à semente aleatória, embora a análise com múltiplas sementes tenha reduzido a incerteza sobre a classificação por MAE. Quarto, o HPO \textit{walk-forward} do XGBoost foi limitado a quatro tentativas completas por custo computacional. Quinto, eventos de alta concentração são raros, o que torna métricas de picos instáveis e dependentes de poucos episódios.

Mesmo com essas limitações, o protocolo atual permite uma conclusão mais defensável que a narrativa anterior: no conjunto final de Sapo, a LSTM direta venceu em erro médio; o XGBoost permaneceu uma referência forte baseada em árvores sobre atributos derivados das janelas temporais; e o novo Seq2Seq atenção ficou competitivo, preservando mais amplitude e melhorando o comportamento em cauda, sem superar a LSTM direta no resultado global.
